
\clearpage
\chapter{מסגרת ההגנה: \textenglish{SkillSieve}}

\section{עקרונות העיצוב}

\textenglish{SkillSieve} מבוסס על שלושה עקרונות:

\begin{itemize}
\item \textbf{בדיקה סטטית}: ניתוח מופעים (\textenglish{AST analysis}) של קוד הכישור ללא הרצה.
\item \textbf{הרצה בסביבת ארגז-חול (\textenglish{sandbox})}: ביצוע הכישור בסביבה מבודדת עם ניטור גישה לרשת ולקבצים.
\item \textbf{ניתוח סמנטי}: שימוש ב-\textenglish{embedding similarity} לזיהוי תיאורים כישורים חריגים.
\end{itemize}

\section{אלגוריתם הסיוג}

ה-\textenglish{SkillSieve} מחשב ציון סיכון לכל כישור:

\begin{equation}
\text{RiskScore}(s) = \alpha \cdot r_{\text{static}}(s) + \beta \cdot r_{\text{sandbox}}(s) + \gamma \cdot r_{\text{semantic}}(s)
\end{equation}

כאשר \(\alpha = 0.4\), \(\beta = 0.4\), \(\gamma = 0.2\) כוילו על מערך בדיקה מאורחת. כישורים עם \(\text{RiskScore} > \theta_{\text{reject}} = 0.65\) נחסמים אוטומטית.

\section{אינטגרציה עם \textenglish{CrewAI}}

ה-\textenglish{SkillSieve} מתממשק לצינור \textenglish{CrewAI} דרך \textenglish{hook} לפני כל קריאה לכישור חיצוני. כל החלטת סיווג נרשמת ב-\textenglish{audit log} לצורך בדיקת ביקורת (\textenglish{compliance audit}) \cite{weidinger2021ethical}.

\begin{figure}[H]
\centering
\includegraphics[width=0.85\textwidth]{assets/architecture.png}
\caption{\textenglish{ClawHavoc Attack Surface and SkillSieve Defense Pipeline}}
\label{fig:arch}
\end{figure}
